\documentclass[../main.tex]{subfiles}
\begin{document}

\appendix
% \appendix does not reliably reset the counter from inside a subfile
\setcounter{section}{0}
\section{Deferred proofs}
\label{app:proofs}

This appendix gives in full the proofs that are compressed or omitted in the main
text. Throughout, $\pi_B$ is the bot prevalence, $c = \pi_B/(1-\pi_B)$,
$\eta = c\rho$, $\lambda = 1-(1+c)\rho$, features are $\mathbf{x}_v \sim
\mathcal{N}(\boldsymbol\mu_{y_v}, \sigma^2 I_d)$ with
$\Delta = \boldsymbol\mu_B - \boldsymbol\mu_H$, and $D$ is the neighbourhood size.

\subsection{The edge-balance identity}
\label{app:balance}

\begin{lemma}
Let camouflage edges be those joining a bot to a human, and suppose each bot directs
a fraction $\rho$ of its $D$ edges to humans. Then the expected fraction of a
human's neighbourhood that is composed of bots is $\eta = c\rho$.
\end{lemma}

\begin{proof}
Let $n$ be the number of nodes, so there are $\pi_B n$ bots and $(1-\pi_B)n$ humans.
Each bot contributes $\rho D$ camouflage edge-endpoints, giving
$\pi_B n \rho D$ endpoints in total. Every such endpoint lands on a human, and the
humans collectively expose $(1-\pi_B) n D$ endpoints. Hence the expected bot fraction
of a human neighbourhood is
\[
    \eta \;=\; \frac{\pi_B n \rho D}{(1-\pi_B) n D}
        \;=\; \frac{\pi_B}{1-\pi_B}\,\rho \;=\; c\rho .
\]
The identity is a conservation statement about edge endpoints and is independent of
how bots choose \emph{which} humans to attach to; only the aggregate count matters.
It presumes $\eta \leq 1$, i.e.\ $\rho \leq 1/c$, which is automatic when
$\pi_B \leq 1/2$.
\end{proof}

The asymmetry is the point: the adversary sets $\rho$ directly but induces $\eta$ on
the other class, and the two coincide only when the classes are balanced ($c=1$).

\subsection{Proof of the depth recursion (Theorem~\ref{thm:depth})}
\label{app:depth}

\begin{proof}
Work with the projection onto the discriminant direction
$\hat{\mathbf{u}} = \Delta/\|\Delta\|$, which is one-dimensional and loses no
information about the class means. Write $m^{(l)}_B, m^{(l)}_H$ for the projected
class-conditional means after $l$ layers and $V^{(l)}_B, V^{(l)}_H$ for the
corresponding variances, with $m^{(0)}_B - m^{(0)}_H = \|\Delta\|$ and
$V^{(0)}_B = V^{(0)}_H = \sigma^2$.

\emph{Means.} A bot aggregates $D$ neighbours of which a fraction $1-\rho$ are bots
and $\rho$ humans, so by linearity
$m^{(l+1)}_B = (1-\rho)m^{(l)}_B + \rho\,m^{(l)}_H$, and symmetrically
$m^{(l+1)}_H = \eta\,m^{(l)}_B + (1-\eta)m^{(l)}_H$. Subtracting,
\[
  \Delta_{l+1} = m^{(l+1)}_B - m^{(l+1)}_H
  = \bigl(1-\rho-\eta\bigr)\bigl(m^{(l)}_B - m^{(l)}_H\bigr)
  = \lambda\,\Delta_l ,
\]
which is~\eqref{eq:depth-mean}; hence $\Delta_L = \lambda^L\|\Delta\|$. Note the two
classes are deflated by different factors ($\rho$ versus $\eta$), so the midpoint
$\tfrac12(m_B + m_H)$ moves unless $c = 1$; the \emph{separation}, however, obeys a
single scalar recursion.

\emph{Variances.} Fix a bot at level $l+1$ and let $X$ be the projected
representation of one of its neighbours. Under the tree-like assumption the $D$
neighbours are independent, and $X$ is a two-component mixture: with probability
$1-\rho$ it has mean $m^{(l)}_B$ and variance $V^{(l)}_B$, and with probability
$\rho$ mean $m^{(l)}_H$ and variance $V^{(l)}_H$. For any such mixture,
\[
  \operatorname{Var}(X)
  = (1-\rho)V^{(l)}_B + \rho V^{(l)}_H
    + \rho(1-\rho)\bigl(m^{(l)}_B - m^{(l)}_H\bigr)^2 ,
\]
by the law of total variance, the last term being the between-component
contribution. Averaging $D$ independent copies divides this by $D$, which
gives~\eqref{eq:depth-varB}; replacing $\rho$ by $\eta$
gives~\eqref{eq:depth-varH}.

\emph{Noise-dominated limit.} Suppose the mixture terms
$\rho(1-\rho)\Delta_l^2$ and $\eta(1-\eta)\Delta_l^2$ are negligible beside the
variance terms at every level. Then both recursions reduce to
$V^{(l+1)} = V^{(l)}/D$, so $V^{(L)} = \sigma^2/D^{L}$ and the standard deviation is
$\sigma D^{-L/2}$. With $\mathrm{SNR}_L := \Delta_L/(2\sqrt{V^{(L)}})$ and
$\mathrm{SNR}_{\mathrm{raw}} = \|\Delta\|/(2\sigma)$,
\[
  \frac{\mathrm{SNR}_L}{\mathrm{SNR}_{\mathrm{raw}}}
  = \frac{\lambda^L\|\Delta\|}{2\sigma D^{-L/2}}\cdot\frac{2\sigma}{\|\Delta\|}
  = \bigl(\lambda\sqrt{D}\bigr)^{L},
\]
which is~\eqref{eq:depth-clean}. Since $t \mapsto t^L$ is increasing on $t>0$, this
exceeds $1$ for some $L \geq 1$ if and only if it does for all of them, so the
crossover $|\lambda|\sqrt{D} = 1$ is independent of $L$; solving
$|1-(1+c)\rho| = D^{-1/2}$ for $\rho$ gives the two endpoints
of~\eqref{eq:window}.
\end{proof}

\begin{remark}[Where the approximation is used, and what it costs]
The tree-like assumption enters only in asserting that the $D$ aggregated
neighbours are independent. On a graph with overlapping receptive fields the
level-$l$ representations of two neighbours share ancestors and are positively
correlated, which inflates $V^{(l+1)}$ above~\eqref{eq:depth-varB} and so makes
\eqref{eq:depth-clean} an \emph{optimistic} bound at large $L$. The recursion is
therefore best read as exact for locally tree-like graphs and as an upper bound on
the achievable SNR otherwise. In particular the unbounded growth implied by
$(\lambda\sqrt D)^L$ for $\lambda\sqrt D>1$ cannot persist on a finite graph, where
receptive fields saturate once they cover it; the depth-independence of the
\emph{crossover}, which is what we use, does not rely on the large-$L$ behaviour.
\end{remark}

\subsection{Exact law of the edge score}
\label{app:score}

\begin{lemma}
For an edge $(u,v)$ let $s_{uv} = \|\mathbf{x}_u - \mathbf{x}_v\|^2$. Then
$s \sim 2\sigma^2\chi^2_d$ if $y_u = y_v$, and
$s \sim 2\sigma^2\chi^2_d(r^2/2)$ with $r = \|\Delta\|/\sigma$ otherwise. In
particular $\mathbb{E}[s \mid \text{camouflage}] - \mathbb{E}[s \mid \text{within}]
= \|\Delta\|^2$, and the ratio of the gap to the within-class standard deviation is
$\kappa = r^2/(2\sqrt{2d})$.
\end{lemma}

\begin{proof}
Write $\mathbf{z} = \mathbf{x}_u - \mathbf{x}_v$. If $y_u = y_v$ then
$\mathbf{z} \sim \mathcal{N}(0, 2\sigma^2 I_d)$, so
$\|\mathbf{z}\|^2/(2\sigma^2) \sim \chi^2_d$. If $y_u \neq y_v$ then
$\mathbf{z} \sim \mathcal{N}(\pm\Delta, 2\sigma^2 I_d)$ and
$\|\mathbf{z}\|^2/(2\sigma^2)$ is noncentral chi-square with $d$ degrees of freedom
and noncentrality $\|\Delta\|^2/(2\sigma^2) = r^2/2$.

For the moments, $\mathbb{E}[\chi^2_d] = d$ and
$\mathbb{E}[\chi^2_d(\lambda)] = d + \lambda$, so the means are $2\sigma^2 d$ and
$2\sigma^2 d + \|\Delta\|^2$, giving the stated gap. Since
$\operatorname{Var}(\chi^2_d) = 2d$, the within-class variance of $s$ is
$4\sigma^4\cdot 2d$, whose square root is $2\sigma^2\sqrt{2d}$; dividing the gap by
it gives $\kappa = \|\Delta\|^2/(2\sigma^2\sqrt{2d}) = r^2/(2\sqrt{2d})$.

The $d^{-1/2}$ dependence is the concentration phenomenon: the gap between the two
score distributions grows like $\|\Delta\|^2$ while their spread grows like
$\sqrt{d}$, so a fixed feature separation becomes progressively harder to detect
from raw squared distances as the ambient dimension increases.
\end{proof}

\section{Reproducibility}
\label{app:repro}

All closed forms in the paper are checked numerically by the scripts in
\texttt{theory/}: \texttt{csbm\_validation.py} covers
Proposition~\ref{prop:deflation} through Theorem~\ref{thm:sgtr},
\texttt{extensions2.py} covers Proposition~\ref{prop:deff},
Corollary~\ref{cor:ceiling} and Theorem~\ref{thm:depth}, and
\texttt{measure\_rho.py} implements the estimator of
Section~\ref{sec:measuring}. Each writes a JSON record of every figure quoted in
the text. The scripts depend only on \texttt{numpy} and \texttt{scipy}.

\end{document}
